% Section 3: GAN and FOMM.
% Facts come from the six notes under fomm/notes/ and from the checked fact sheet.
% The old slides are not used where the fact check corrected them.
% Figure inputs are relative to docs/report/.

\section{GAN \& FOMM}

A generator is like a painter who draws faces. A discriminator is like an expert
who judges each face as real or fake and teaches the painter. This game cannot
follow orders, however. For reenactment we must keep one identity and copy the
motion from a video, so we used the First Order Motion Model (FOMM).

\subsection{Standard GANs}
A standard GAN holds two networks. The generator $G$ reads a noise vector $z$
and draws a face; the discriminator $D$ reads an image and gives a verdict: real
or fake. The two train against each other, so the generator learns to fool the
discriminator (Goodfellow et al., 2014), as Figure~\ref{fig:gan_training} shows.

We cannot control the output. Identity, pose and mood are all mixed into the
noise input. We cannot say ``keep this person, copy that motion''. A
reenactment model needs two inputs instead: a source image (who) and a driving
video (how it moves).

% Figure: the standard GAN training loop (Goodfellow et al., 2014).
% The Generator turns noise into a fake face. The Discriminator compares the fake
% face with real faces from the dataset. The feedback teaches the Generator.
% Requires in the preamble: \usepackage{tikz} + \usetikzlibrary{arrows.meta}
\begin{figure}[tb]
  \centering
  \begin{tikzpicture}[
    x=1cm, y=1cm,
    font=\small,
    block/.style={rectangle, rounded corners=2pt, draw=black!70, line width=0.5pt,
      fill=blue!10, align=center, inner sep=4pt, minimum height=9mm},
    inputblock/.style={block, fill=green!12},
    stageblock/.style={block, fill=orange!15},
    flow/.style={->, >=Latex, line width=0.7pt, draw=black!75},
    feed/.style={->, >=Latex, line width=0.7pt, draw=black!75, dashed},
  ]

  % ---------------- the data ----------------
  \node[inputblock, text width=1.20cm, minimum height=1.05cm] (z) at (0.75,0)
    {\textbf{noise} $z$};
  \node[inputblock, text width=2.05cm, minimum height=1.05cm] (data) at (1.20,2.10)
    {\textbf{Dataset} of\\real faces};
  \node[inputblock, text width=1.45cm, minimum height=1.05cm] (real) at (6.30,2.10)
    {real face};

  % ---------------- the two networks ----------------
  \node[block, text width=1.90cm, minimum height=1.05cm] (G) at (3.55,0)
    {\textbf{Generator}\\$G$};
  \node[block, text width=2.50cm, minimum height=1.05cm] (D) at (9.90,1.05)
    {\textbf{Discriminator}\\$D$};

  % ---------------- the fake face and the verdict ----------------
  \node[inputblock, text width=1.45cm, minimum height=1.05cm] (fake) at (6.30,0)
    {fake face};
  \node[stageblock, text width=2.10cm, minimum height=1.05cm] (verdict) at (13.05,1.05)
    {\textbf{verdict}:\\real or fake};

  % ---------------- arrows ----------------
  \draw[flow] (z) -- (G);
  \draw[flow] (G) -- (fake);
  \draw[flow] (data) -- (real);
  \draw[flow] (fake) -- (D);
  \draw[flow] (real) -- (D);
  \draw[flow] (D) -- (verdict);

  % ---------------- the training feedback ----------------
  \draw[feed] (D.south) -- (9.90,-0.85) -- (3.55,-0.85) -- (G.south);
  \node[font=\scriptsize, align=center, anchor=north] at (6.72,-0.95)
    {the Generator learns to\\fool the Discriminator};

  \end{tikzpicture}
  \caption{The training loop of a standard GAN. The Discriminator compares each
    fake face with real faces and gives a verdict; the dashed feedback teaches
    the Generator to fool it.}
  \label{fig:gan_training}
\end{figure}


\subsection{First Order Motion Model}
FOMM animates one source photo with the head motion and the facial expressions
from a driving video. The method comes from Siarohin et al. (``First Order
Motion Model for Image Animation'', NeurIPS 2019). It needs no landmark detector
and no 3D face model.
It learns its own keypoints per object category; our category is faces. It
follows Monkey-Net, which the paper calls the first object-agnostic deep model,
and FOMM improves on it.

\subsubsection{First-order motion}
0th order motion moves a point only in $x$ and $y$, so the skin keeps its
shape and the result looks stiff and robotic. 1st order motion also measures how
the skin around the point turns, grows and stretches, through a small
$2\times2$ Jacobian matrix per keypoint. The paper calls it a first order Taylor
expansion.

\subsubsection{The three parts}
The pipeline has three parts (Figure~\ref{fig:fomm_pipeline}).

% Figure: the First Order Motion Model pipeline (Siarohin et al., NeurIPS 2019).
% The keypoint detector reads the source image once and each driving frame.
% The dense motion network makes a motion field and an occlusion map.
% The generator warps the source features and inpaints the occluded parts.
% Requires in the preamble: \usepackage{tikz} + \usetikzlibrary{arrows.meta}
\begin{figure}[tb]
  \centering
  \begin{tikzpicture}[
    x=1cm, y=1cm,
    font=\small,
    block/.style={rectangle, rounded corners=2pt, draw=black!70, line width=0.5pt,
      fill=blue!10, align=center, inner sep=4pt, minimum height=9mm},
    inputblock/.style={block, fill=green!12},
    noteblock/.style={rectangle, rounded corners=2pt, draw=black!45, dashed,
      fill=red!8, align=left, inner sep=3pt, font=\scriptsize},
    flow/.style={->, >=Latex, line width=0.7pt, draw=black!75},
  ]

  % ---------------- inputs (left) ----------------
  \node[inputblock, text width=2.00cm, minimum height=1.10cm] (src) at (1.15,1.60)
    {\textbf{Source image}\\[1pt]\scriptsize the identity};
  \node[inputblock, text width=2.00cm, minimum height=1.10cm] (drv) at (1.15,-0.10)
    {\textbf{Driving video}\\[1pt]\scriptsize the motion};

  % ---------------- the three parts ----------------
  \node[block, text width=2.95cm, minimum height=2.30cm] (kpd) at (5.10,0.75)
    {\textbf{Keypoint detector}\\[1pt]\scriptsize U-Net network\\$K=10$ keypoints\\one Jacobian each};
  \node[block, text width=3.05cm, minimum height=2.30cm] (dmn) at (9.10,0.75)
    {\textbf{Dense motion}\\\textbf{network}\\[1pt]\scriptsize motion field (optical\\flow) + occlusion map\\with values in $[0,1]$};
  \node[block, text width=2.35cm, minimum height=1.50cm] (gen) at (12.75,0.75)
    {\textbf{Generator}\\[1pt]\scriptsize warp + inpaint};

  % ---------------- output ----------------
  \node[inputblock, text width=2.10cm, minimum height=1.00cm] (out) at (12.75,-1.60)
    {\textbf{Output frame}};

  % ---------------- arrows ----------------
  \draw[flow] (src.east) -- (kpd.west |- src.east);
  \draw[flow] (drv.east) -- (kpd.west |- drv.east);
  \draw[flow] (kpd) -- (dmn);
  \draw[flow] (dmn) -- (gen);
  \draw[flow] (gen.south) -- (out.north);

  % ---------------- notes ----------------
  \node[noteblock, text width=5.30cm, anchor=north west] at (0.20,-1.30)
    {The detector reads the source once and every driving frame.};
  \node[noteblock, text width=5.20cm, anchor=north west] at (5.90,-1.30)
    {Each frame is processed alone; no memory between frames.};

  \end{tikzpicture}
  \caption{The First Order Motion Model pipeline: keypoints with Jacobians, a
    motion field and an occlusion map, then warp and inpaint. Each frame is
    processed alone.}
  \label{fig:fomm_pipeline}
\end{figure}


The \textbf{keypoint detector} is a U-Net network (the paper's word; the code
names the class Hourglass). It finds $K=10$ keypoints for the faces model and
one Jacobian per keypoint. The expected point position is a weighted average
over the output heatmaps (soft-argmax). An equivariance loss teaches the
detector during training, so it needs no keypoint labels.

The \textbf{dense motion network} turns the sparse keypoints into a dense motion
field (optical flow) and an occlusion map with values between 0 and 1.

The \textbf{generator} warps the source features with the flow. The code
multiplies the warped features by the occlusion map
(\texttt{out = out * occlusion\_map}). Where the map is near 0, the warped
pixels lose their weight and the decoder fills the area from learned face priors
(inpainting).

Each frame is processed on its own. No state passes from frame to frame (no
recurrence).

\subsubsection{Occlusion and inpainting}
A closed mouth or a turned head reveals pixels that do not exist in the source
photo (self-occlusion). Where $O$ is near 1, the pixels come from the source.
Where $O$ is near 0, the network draws the area itself from the training-set
average (Figure~\ref{fig:fomm_occlusion}).

These are the failure modes we saw. Rare personal details are lost in inpainted
areas. Very large head turns, about 45 degrees or more, leave blurry patches.
Both limits come from our own observation, not from the paper.

% Figure: occlusion and inpainting in the FOMM generator.
% A closed mouth or a turned head reveals pixels that are missing in the source.
% $O \approx 1$: the pixels come from the source. $O \approx 0$: the network
% draws the area itself from the training-set average.
% Requires in the preamble: \usepackage{tikz} + \usetikzlibrary{arrows.meta}
\begin{figure}[tb]
  \centering
  \begin{tikzpicture}[
    x=1cm, y=1cm,
    font=\small,
    block/.style={rectangle, rounded corners=2pt, draw=black!70, line width=0.5pt,
      fill=blue!10, align=center, inner sep=4pt, minimum height=9mm},
    inputblock/.style={block, fill=green!12},
    flow/.style={->, >=Latex, line width=0.7pt, draw=black!75},
    leader/.style={draw=black!45, line width=0.4pt},
  ]

  % ---------------- source frame (closed mouth) ----------------
  \node[inputblock, minimum width=2.20cm, minimum height=2.20cm] (src) at (1.15,0.90) {};
  \node[font=\scriptsize] at (1.15,1.72) {\textbf{Source image}};
  \fill[black!70] (0.80,1.22) circle (0.6mm);
  \fill[black!70] (1.50,1.22) circle (0.6mm);
  \draw[black!70, line width=1pt] (0.80,0.42) -- (1.50,0.42);
  \node[font=\scriptsize] at (1.15,0.10) {closed mouth};

  % ---------------- driving frame (open mouth) ----------------
  \node[inputblock, minimum width=2.20cm, minimum height=2.20cm] (drv) at (1.15,-1.60) {};
  \node[font=\scriptsize] at (1.15,-0.72) {\textbf{Driving frame}};
  \fill[black!70] (0.80,-1.22) circle (0.6mm);
  \fill[black!70] (1.50,-1.22) circle (0.6mm);
  \fill[draw=black!60, line width=0.4pt, orange!35] (1.15,-2.02) ellipse (0.42cm and 0.26cm);
  \node[font=\scriptsize] at (1.15,-2.47) {open mouth};

  % ---------------- generator ----------------
  \node[block, text width=2.20cm, minimum height=1.40cm] (gen) at (4.55,-0.35)
    {\textbf{Generator}\\[1pt]\scriptsize warp + inpaint};
  \draw[flow] (src.east) -| (gen.north);
  \draw[flow] (drv.east) -| (gen.south);

  % ---------------- output frame ----------------
  \node[inputblock, minimum width=2.60cm, minimum height=2.70cm] (out) at (7.75,-0.35) {};
  \node[font=\scriptsize] at (7.75,0.62) {\textbf{Output frame}};
  \fill[black!70] (7.35,0.12) circle (0.6mm);
  \fill[black!70] (8.15,0.12) circle (0.6mm);
  \fill[draw=black!60, line width=0.4pt, orange!35] (7.75,-0.80) ellipse (0.55cm and 0.33cm);
  \draw[flow] (gen.east) -- (out.west);

  % ---------------- the two regions ----------------
  \draw[leader] (7.75,1.00) -- (7.75,1.14);
  \node[font=\scriptsize, align=center, text width=4.30cm, anchor=south]
    at (7.75,1.16) {visible pixels: taken from the source, $O\approx1$};
  \draw[leader] (8.30,-0.80) -- (9.45,-0.80);
  \node[font=\scriptsize, align=left, text width=3.30cm, anchor=west]
    at (9.50,-0.80) {new area: no pixels in the source, $O\approx0$, drawn by the network};

  \end{tikzpicture}
  \caption{Occlusion and inpainting. Where $O$ is near 1, the pixels come from
    the source. Where $O$ is near 0, the network draws the missing pixels
    itself.}
  \label{fig:fomm_occlusion}
\end{figure}


\subsection{Local experiment}
We ran the released VoxCeleb checkpoint \texttt{vox-cpk.pth.tar} in CPU mode on
an Apple Silicon machine. The source portrait is 512x512; the driving clip is
256x256, 320 frames, 29.97 FPS. With \texttt{-{}-relative} the keypoints move
relative to the first frame; \texttt{-{}-adapt\_scale} adjusts the motion size
with the convex hull of the keypoints.

The output is 256x256, 320 frames, 29.97 FPS and 10.7 seconds long. The recorded
speed is about 2.05 iterations per second, about 2.5 minutes in total. Head
turns and blinks transferred smoothly; mouth motion needed inpainting.
Table~\ref{tab:fomm_run} lists the setup and the output.

\begin{table}[tb]
  \centering
  \footnotesize
  \begin{tabular}{@{}l l@{}}
    \toprule
    Item & Value \\
    \midrule
    Checkpoint & \texttt{vox-cpk.pth.tar} (VoxCeleb) \\
    Source portrait & 512 $\times$ 512 \\
    Driving clip & 256 $\times$ 256, 320 frames, 29.97 FPS \\
    Mode and flags & CPU, \texttt{-{}-relative}, \texttt{-{}-adapt\_scale} \\
    Output & 256 $\times$ 256, 320 frames, 29.97 FPS, 10.7 s \\
    Speed & about 2.05 it/s, about 2.5 minutes \\
    \bottomrule
  \end{tabular}
  \caption{The local FOMM run: setup, flags and output.}
  \label{tab:fomm_run}
\end{table}

\subsection{Failure mode: braces}
The VoxCeleb training data has no braces, so the mouth area should lose them
when inpainting starts. Our frames confirm this (Table~\ref{tab:fomm_braces}).

\begin{table}[tb]
  \centering
  \footnotesize
  \begin{tabular}{@{}r r l l@{}}
    \toprule
    Frame & Time & Driving mouth & Active path and result \\
    \midrule
    10 & 0.3 s & small shift & warping, $O\approx1$: braces kept \\
    195 & 6.5 s & lips close & mouth hidden \\
    288 & 9.6 s & wide smile & inpainting, $O\approx0$: plain teeth \\
    \bottomrule
  \end{tabular}
  \caption{The braces test: the active path decides where the mouth pixels come
    from.}
  \label{tab:fomm_braces}
\end{table}

The active path decides which pixels appear: warped areas keep the pixels of the
person, and inpainted areas follow the training data. FOMM gave the project a
fast per-frame method. Identity details survive only in warped areas, and rare
details such as braces are lost as soon as inpainting starts. There is no memory
across frames, so the method cannot repair a frame from its neighbours.
